%! TEX root = main.tex

\subsection{Motivation: what problem is machine learning solving?}

Everything that follows is an answer to one question. We can measure a
\emph{loss}: a single number that says how badly some function $f_\theta$
performs on data. How do we find the parameters $\theta$ that make it small?

This reframes what you already know. ``A neural network is
$w\cdot x+b$'' describes one \emph{function} in a large \emph{family} of
functions. Machine learning is not about that function; it is about the
optimisation problem sitting on top of it. Three ingredients are needed:
\begin{enumerate}[label=\arabic*.]
  \item a family of candidate functions, called the \emph{model} or
    \emph{hypothesis class}, indexed by parameters $\theta$;
  \item a \emph{loss} that turns ``how wrong is $f_\theta$ on this example''
    into a real number;
  \item an algorithm that searches $\theta$-space for a low loss.
\end{enumerate}
Matrices enter in step~1, because a linear layer evaluated on a batch of
inputs is a matrix product. Backpropagation enters in step~3, because you
cannot search $\R^{d}$ blindly when $d$ is $10^9$: you need the gradient,
and backprop is the algorithm that computes it.

The roadmap for this topic:
\begin{enumerate}[label=\arabic*.]
  \item supervised learning as loss minimisation, and the linear layer as a
    matrix (this lesson);
  \item the chain rule in matrix form, and backpropagation as reverse-mode
    automatic differentiation (this lesson);
  \item gradient descent, minibatching, and then generalisation---why small
    training loss is not the goal;
  \item the same machinery for sequential decisions: dynamic programming,
    $Q$-learning, and the policy gradient, which is backprop applied to a
    stochastic objective.
\end{enumerate}

\subsection{The supervised learning contract}

\begin{definition}[Supervised learning problem]
An \emph{example} is a pair $(\mathbf{x}, \mathbf{y})$ with
$\mathbf{x}\in\R^{n}$ (the input or features) and $\mathbf{y}\in\R^{k}$
(the target). A \emph{dataset} is a finite list
$D=((\mathbf{x}_1,\mathbf{y}_1),\ldots,(\mathbf{x}_m,\mathbf{y}_m))$.

Fix a \emph{model} $f_\theta\colon\R^{n}\to\R^{k}$ depending on parameters
$\theta\in\R^{d}$, and a \emph{per-example loss}
$\ell\colon\R^{k}\times\R^{k}\to\R_{\ge 0}$. The \emph{empirical risk} of
$\theta$ is
\[
  L_D(\theta)=\frac{1}{m}\sum_{i=1}^{m}\ell\bigl(f_\theta(\mathbf{x}_i),\mathbf{y}_i\bigr).
\]
\emph{Training} means solving, at least approximately,
$\min_{\theta}L_D(\theta)$.
\end{definition}

Two comments on this definition, both of which matter later.

First, $L_D$ is an \emph{average}, not a sum. This is only a choice of scale,
but it means the loss does not blow up when you change the batch size, and an
``average gradient'' is genuinely an average.

Second, $L_D$ is a function of one variable: the flat vector $\theta$. All the
interesting structure (matrices, layers, activations) is hidden inside
$f_\theta$. This is why the mathematics of training is \emph{calculus plus
linear algebra}, and why we will never need a separate theory for each
architecture: only $\partial L/\partial\theta$ is needed, however $f_\theta$
is built.

\begin{example}[Losses you will meet]
\leavevmode
\begin{enumerate}[label=(\alph*)]
  \item \emph{Squared error}, for regression: $\ell(\hat{\mathbf{y}},\mathbf{y})=\tfrac12\norm{\hat{\mathbf{y}}-\mathbf{y}}^2$.
    The factor $\tfrac12$ cancels the $2$ from differentiating a square; it is
    pure convenience and changes no minimiser.
  \item \emph{Mean squared error} (MSE) over a batch: $L=\frac{1}{2m}\norm{Z-Y}_F^2$,
    where the rows of $Z\in\R^{m\times k}$ are the predictions and the rows
    of $Y$ are the targets, and $\norm{\cdot}_F$ is the Frobenius norm
    $\norm{A}_F^2=\sum_{ij}A_{ij}^2$.
  \item \emph{Softmax cross-entropy}, for $C$-way classification, which we
    will build in a later lesson from the same ingredients.
\end{enumerate}
\end{example}

The Frobenius norm is worth noticing: it is the Euclidean norm of the matrix
flattened into a vector. So squaring it is exactly ``sum of squares over all
entries'', and MSE is ``average squared coordinate error''.

\subsection{A layer is a matrix, a batch is a stacked matrix}

\begin{definition}[Affine (fully connected) layer]
An affine layer with $n$ inputs and $p$ outputs has parameters
$W\in\R^{n\times p}$ (the \emph{weight matrix}) and $\mathbf{b}\in\R^{p}$
(the \emph{bias}). It maps a single input $\mathbf{x}\in\R^{n}$ to
\[
  \mathbf{z}=W^{\!\top}\mathbf{x}+\mathbf{b}\in\R^{p}.
\]
\end{definition}

Note the transposition: with $\mathbf{x}$ as a column, the $j$-th output
coordinate is $z_j=\sum_{k=1}^{n}W_{kj}x_k+b_j$, i.e.\ the inner product of
$\mathbf{x}$ with the $j$-th \emph{column} of $W$. In your linear algebra
course the layer would naturally be written as $A\mathbf{x}+\mathbf{b}$ with
$A\in\R^{p\times n}$; the convention above stores $A^{\!\top}$ instead so
that the batch version needs no transposition. Both are the same map; only the
storage differs.

Now feed a whole batch at once. Put the inputs as the \emph{rows} of
$X\in\R^{m\times n}$ and the outputs as the rows of $Z\in\R^{m\times p}$. Then
\[
  Z=XW+\mathbf{1}_m\mathbf{b}^{\!\top},
\]
where $\mathbf{1}_m\in\R^{m}$ is the all-ones vector, so
$\mathbf{1}_m\mathbf{b}^{\!\top}\in\R^{m\times p}$ has every row equal to
$\mathbf{b}^{\!\top}$. Row $i$ of this identity reads
$\mathbf{z}_i^{\!\top}=\mathbf{x}_i^{\!\top}W+\mathbf{b}^{\!\top}$, which is
the single-example layer transposed. This is the whole content of the next
proposition.

\begin{proposition}[Batching is not a new operation]\label{prop:batching}
Let $X\in\R^{m\times n}$ have rows $\mathbf{x}_1^{\!\top},\ldots,\mathbf{x}_m^{\!\top}$.
Then the rows of $Z=XW+\mathbf{1}_m\mathbf{b}^{\!\top}$ are exactly
$f(\mathbf{x}_1)^{\!\top},\ldots,f(\mathbf{x}_m)^{\!\top}$ for the affine layer
$f$ above. Consequently the batch computation is one matrix multiplication
whose arithmetic cost is $\Theta(mnp)$, i.e.\ exactly $m$ times the cost of a
single example, with no saving from sharing the work.
\end{proposition}

\begin{proof}
Row $i$ of $XW$ is $\sum_{k=1}^{n}X_{ik}W_{k\cdot}=\mathbf{x}_i^{\!\top}W$, and
row $i$ of $\mathbf{1}_m\mathbf{b}^{\!\top}$ is $\mathbf{b}^{\!\top}$. Add them.
Independence across $i$ holds because the $i$-th row of $Z$ depends only on the
$i$-th row of $X$.
\end{proof}

\begin{remark}[Why rows and not columns]
Storing examples as rows is a convention, and PyTorch, JAX and TensorFlow all
follow it. The concrete benefit is that the feature axis $n$ is the
\emph{second} index of $X$, so the $n$ entries belonging to one example are
contiguous in a row-major array, and the whole layer is a single BLAS call
$XW$ with no transposition of the batch. The cost is conceptual: the elementary
linear-algebra object ``apply the map to $\mathbf{x}$'' becomes ``multiply the
row by $W$'', which is why the transpose in the definition of the layer cannot
be avoided. Pick one convention, write it down, and check shapes against it;
mixing the two is the source of most dimension errors.
\end{remark}

\begin{remark}[Why a nonlinearity is mandatory]
By \cref{prop:batching} a layer is affine. Compose two affine layers:
\[
  \mathbf{z}=W_2^{\!\top}\bigl(W_1^{\!\top}\mathbf{x}+\mathbf{b}_1\bigr)+\mathbf{b}_2
  =(W_1W_2)^{\!\top}\mathbf{x}+\bigl(W_2^{\!\top}\mathbf{b}_1+\mathbf{b}_2\bigr).
\]
The composition is again affine, with weight $W_1W_2$. So depth without
nonlinearity buys nothing: you are still fitting a single affine map, now with
a rank constraint ($\operatorname{rank}(W_1W_2)\le\min(\operatorname{rank}W_1,\operatorname{rank}W_2)$).
This is the same ``rank cannot increase under composition'' argument as in your
linear algebra exercises, and it is the reason every real network interleaves
a nonlinear map.
\end{remark}

\begin{definition}[Activation, and an MLP]
An \emph{activation} is a nonlinear scalar function $\sigma$ applied entrywise
to a matrix, written $\sigma(Z)$ with $\sigma(Z)_{ij}=\sigma(Z_{ij})$. The
standard choice here is the rectified linear unit
\[
  \operatorname{relu}(t)=\max(t,0).
\]
A \emph{multi-layer perceptron} (MLP) with one hidden layer is
\begin{align*}
  Z_1&=XW_1+\mathbf{1}\mathbf{b}_1^{\!\top}, &
  R&=\operatorname{relu}(Z_1), &
  Z_2&=RW_2+\mathbf{1}\mathbf{b}_2^{\!\top},
\end{align*}
with $W_1\in\R^{n\times h}$, $W_2\in\R^{h\times k}$. $Z_1$ is called the
\emph{pre-activation} and $R$ the \emph{hidden representation}.
\end{definition}

Why does this help, given that each layer is affine? Entrywise $\operatorname{relu}$
applied to a matmul output is a \emph{piecewise}-affine map. Fix the pattern of
which entries are positive (a boolean matrix, the \emph{mask}); on the region
of inputs $X$ where that pattern is constant, the network reduces to a single
affine map, and different regions give different affine maps. So an MLP is an
affine spline: many linear pieces glued together. The mask is a piece of data
depending on the input, which is precisely what a fixed matrix cannot express.

\begin{remark}[Parameter counting]
The first layer has $nh+h$ parameters, the second $hk+k$, total
$nh+h+hk+k$. For $n=784$, $h=128$, $k=10$ this is about $10^5$; for language
models the same formula with $h\sim10^{4}$ in dozens of layers yields
$10^{9}$--$10^{11}$. Number the parameters in some fixed order and stack them
into $\theta\in\R^{d}$. Gradients will have shape $d$ too, and the whole point
of backprop is to get all $d$ of them for roughly the cost of one forward pass.
\end{remark}

\subsection{The objective is a scalar, so the gradient is a vector}

\begin{definition}[Gradient]
For a differentiable $L\colon\R^{d}\to\R$, the gradient $\nabla L(\theta)$ is
the vector in $\R^{d}$ whose $j$-th entry is
$\partial L/\partial\theta_j$. The \emph{total derivative} at $\theta$ is the
row vector $\nabla L(\theta)^{\!\top}$, the unique linear map
$\R^{d}\to\R$ with
\[
  L(\theta+\delta)=L(\theta)+\nabla L(\theta)^{\!\top}\delta+o(\norm{\delta}).
\]
\end{definition}

Everything downstream depends on one structural fact: \emph{the loss is
scalar}. Whatever the shape of an intermediate quantity, the derivative of
$L$ with respect to it has the same shape as the quantity. This is why we can
avoid the $4$-index tensors that ``derivative of a matrix with respect to a
matrix'' would otherwise require.

\begin{definition}[Adjoint]
Let $L$ be the scalar loss and let $G$ be any intermediate matrix or vector in
the computation. Write
\[
  \bar{G}=\frac{\partial L}{\partial G},\qquad \bar{G}_{ij}=\frac{\partial L}{\partial G_{ij}},
\]
so $\bar{G}$ has the \emph{same shape as} $G$. We call $\bar{G}$ the
\emph{adjoint} of $G$, also called its \emph{upstream gradient}. For the
parameters, $\bar{W}$ and $\bar{\mathbf b}$ are the gradient pieces we want.
\end{definition}

The adjoint is not the naive ``derivative matrix''. If $G\in\R^{m\times p}$
then $G$ has $mp$ entries; a linear map $\R^{mp}\to\R$ is described by $mp$
numbers, and $\bar G$ is exactly that description. In linear algebra language,
$\bar G$ is the element of the dual space $(\R^{m\times p})^*$ corresponding to
$\mathrm{d}L$, transported back to $G$'s own shape. That identification is what
makes ``same shape'' a legal convention.

\begin{lemma}[Chain rule, adjoint form]\label{lem:adjoint-chain}
Suppose $L$ depends on $G$ and $G$ depends on $H$ through a differentiable map
$\Phi$, so that $G=\Phi(H)$. Then the adjoints satisfy
\[
  \bar{H}=\bigl(D\Phi(H)\bigr)^{\!\top}\bar{G},
\]
where $D\Phi(H)$ is the Jacobian of $\Phi$ at $H$ (a linear map between the
flattened spaces) and $^{\!\top}$ denotes its transpose (the adjoint map).
\end{lemma}

\begin{proof}
Flatten $H$ and $G$ into vectors $\mathbf{h},\mathbf{g}$. By definition
$D\Phi$ is the matrix with $\partial g_a/\partial h_b$ in position $(a,b)$.
Then for each $b$,
\[
  \frac{\partial L}{\partial h_b}
  =\sum_{a}\frac{\partial L}{\partial g_a}\frac{\partial g_a}{\partial h_b}
  =\sum_{a}\bigl(D\Phi\bigr)_{ab}\,\bar{G}_{a}
  =\bigl((D\Phi)^{\!\top}\bar{G}\bigr)_{b}. \qedhere
\]
\end{proof}

So backpropagation is one line of mathematics: \emph{multiply the upstream
adjoint by the transpose of the Jacobian, at every step, moving backwards}.
The only work is computing $D\Phi^{\!\top}\bar G$ for each kind of layer without
ever forming $D\Phi$. The saving is enormous. For a layer with $n=10^{3}$
inputs, $p=10^{3}$ outputs and a batch of $m=10^{4}$ examples, the Jacobian of
$W\mapsto Z$ is the $(mp)\times(np)$ matrix of size $10^{7}\times10^{6}$, while
the product $X^{\!\top}\bar Z$ is an $n\times p$ matrix computed with one
matmul. Reverse mode never materialises the Jacobian; it only ever applies its
transpose to a vector.

\begin{theorem}[Backprop through an affine layer]\label{thm:affine-backprop}
Let $Z=XW+\mathbf{1}_m\mathbf{b}^{\!\top}$ with $X\in\R^{m\times n}$,
$W\in\R^{n\times p}$, $\mathbf{b}\in\R^{p}$, and let $\bar Z$ be the adjoint of
$Z$. Then
\[
  \bar W=X^{\!\top}\bar Z,\qquad
  \bar{\mathbf b}=\bar Z^{\!\top}\mathbf{1}_m,\qquad
  \bar X=\bar Z\,W^{\!\top}.
\]
\end{theorem}

\begin{proof}
Since $Z_{ij}=\sum_{k}X_{ik}W_{kj}+b_j$, we have
$\partial Z_{ij}/\partial W_{kj'}=\delta_{jj'}X_{ik}$,
$\partial Z_{ij}/\partial b_{j'}=\delta_{jj'}$, and
$\partial Z_{ij}/\partial X_{ik'}=\delta_{kk'}W_{kj}$. Summing over all
$(i,j)$ with the adjoint weights $\bar Z_{ij}$:
\begin{align*}
  \bar W_{kj}
  &=\sum_{i}\sum_{j'}\bar Z_{ij'}\frac{\partial Z_{ij'}}{\partial W_{kj}}
   =\sum_{i}\bar Z_{ij}X_{ik}=(X^{\!\top}\bar Z)_{kj},\\
  \bar b_{j}
  &=\sum_{i}\sum_{j'}\bar Z_{ij'}\frac{\partial Z_{ij'}}{\partial b_{j}}
   =\sum_{i}\bar Z_{ij}=(\bar Z^{\!\top}\mathbf{1}_m)_{j},\\
  \bar X_{ik}
  &=\sum_{i'}\sum_{j}\bar Z_{i'j}\frac{\partial Z_{i'j}}{\partial X_{ik}}
   =\sum_{j}\bar Z_{ij}W_{kj}=(\bar Z W^{\!\top})_{ik}. \qedhere
\end{align*}
\end{proof}

\begin{remark}[Reading the three formulas]
All three are ``sum over the index the layer contracted away''.
\begin{itemize}
  \item $\bar W=X^{\!\top}\bar Z$: sum the outer products $\mathbf{x}_i\bar{\mathbf{z}}_i^{\!\top}$
    over the batch. Each example contributes the change it would like in $W$.
  \item $\bar{\mathbf b}=\bar Z^{\!\top}\mathbf{1}_m$: column sums of $\bar Z$. The
    bias is added to every example identically, so its gradient is the total
    over the batch.
  \item $\bar X=\bar ZW^{\!\top}$: the same map as the forward pass but with $W$
    transposed, pushing the error back to the inputs. If you do not need
    $\bar X$ (the first layer), skip it---this is the one step you can omit.
\end{itemize}
Shapes are a complete correctness check: $X^{\!\top}\bar Z$ is
$(n\times m)(m\times p)=n\times p$, matching $W$; $\bar ZW^{\!\top}$ is
$(m\times p)(p\times n)=m\times n$, matching $X$. A dimension mismatch always
means you transposed the wrong factor.
\end{remark}

\begin{theorem}[Backprop through an entrywise activation]\label{thm:relu-backprop}
Let $R=\sigma(Z)$ with $\sigma$ applied entrywise and differentiable. Then
\[
  \bar Z=\bar R\odot\sigma'(Z),
\]
where $\odot$ is the Hadamard (entrywise) product and $\sigma'(Z)$ is the matrix
of entrywise derivatives. For $\sigma=\operatorname{relu}$,
$\sigma'(t)=\mathbf{1}[t>0]$, so
\[
  \bar Z_{ij}=\bar R_{ij}\cdot\mathbf{1}[Z_{ij}>0].
\]
\end{theorem}

\begin{proof}
Since $R_{ij}$ depends only on $Z_{ij}$, and $R_{ij'}$ for $j'\ne j$ does not
depend on $Z_{ij}$, the Jacobian is diagonal within each entry:
$\partial R_{ij}/\partial Z_{ij}=\sigma'(Z_{ij})$ and
$\partial R_{ij}/\partial Z_{i'j'}=0$ if $(i,j)\ne(i',j')$. Hence
$\bar Z_{ij}=\sum_{i'j'}\bar R_{i'j'}\partial R_{i'j'}/\partial Z_{ij}
=\bar R_{ij}\sigma'(Z_{ij})$.
\end{proof}

\begin{remark}[The ReLU derivative at $0$]
The function $\operatorname{relu}$ is not differentiable at $0$, and no
derivative value there is ``correct''. Every framework picks $0$, which means
an exactly-zero pre-activation receives no gradient. This is a convention, not a
theorem, and it is the usual way a unit dies (\cref{ex:ml-6}). All our formulas
use $\mathbf{1}[t>0]$, matching that convention.
\end{remark}

The ReLU adjoint is a mask multiply, and that is the whole story of
``nonlinearity'' in the backward pass: the forward pass decides which units are
on, and the backward pass lets gradient through exactly those.

\begin{example}[Adjunction is not optional: a wrong-but-plausible formula]
A common mistake is to write $\bar Z=\sigma'(\bar R)$ or to forget that the
mask must use the \emph{forward} values $Z$, not $\bar R$. \Cref{thm:relu-backprop}
shows both are wrong: the mask is a fixed $0/1$ pattern determined in the
forward pass, and it multiplies the incoming gradient entrywise. Similarly,
writing $W_2\bar R$ instead of $\bar R W_2^{\!\top}$ has incompatible shapes
(unless $m=k$, and then it is silently wrong). \emph{Check shapes before you
check algebra.}
\end{example}

\subsection{The full backward pass of a one-hidden-layer MLP}

Assemble the pieces. With $Z_1=XW_1+\mathbf{1}\mathbf{b}_1^{\!\top}$,
$R=\operatorname{relu}(Z_1)$, $Z_2=RW_2+\mathbf{1}\mathbf{b}_2^{\!\top}$, and the
mean squared error objective
\[
  L=\frac{1}{2m}\norm{Z_2-Y}_F^2,
\]
the adjoints are, in reverse order:
\begin{align}
  \bar Z_2 &= \frac{1}{m}\bigl(Z_2-Y\bigr),
  \label{eq:mse-adjoint}\\
  \bar W_2 &= R^{\!\top}\bar Z_2, \qquad \bar{\mathbf b}_2=\bar Z_2^{\!\top}\mathbf{1}_m, \label{eq:w2}\\
  \bar R &= \bar Z_2 W_2^{\!\top}, \label{eq:Rbar}\\
  \bar Z_1 &= \bar R\odot\mathbf{1}[Z_1>0], \label{eq:z1bar}\\
  \bar W_1 &= X^{\!\top}\bar Z_1, \qquad \bar{\mathbf b}_1=\bar Z_1^{\!\top}\mathbf{1}_m. \label{eq:w1}
\end{align}
The only ad hoc line is \eqref{eq:mse-adjoint}; it is a direct computation,
which we do in the exercises. Everything else is
\cref{thm:affine-backprop} and \cref{thm:relu-backprop}.

\begin{remark}[Why this is fast]
Let $d$ be the number of parameters and $T_{\text{fwd}}$ the cost of evaluating
$L$. The backward pass costs about $2T_{\text{fwd}}$: each layer theorem does
one matrix multiply whose cost is comparable to the forward one. Crucially the
cost is \emph{independent of $d$} as a separate factor---it is proportional to
the size of the computation, which already scales with $d$. A naive
finite-difference gradient estimate, by contrast, needs one full forward pass
per parameter: $O(d)\cdot T_{\text{fwd}}$.

The structural reason is that $L$ has one output. Forward-mode differentiation
propagates a directional derivative $\langle\nabla L,\mathbf{v}\rangle$ (a
\emph{Jacobian-vector product}), which is cheap per input direction but must be
repeated for each of the $d$ parameters. Reverse-mode propagates
$D\Phi^{\!\top}\bar G$ (a \emph{vector-Jacobian product}), which is cheap per
output. Since the output is one scalar, one reverse sweep suffices. This is the
entire content of ``backpropagation'', and it is exactly reverse-mode automatic
differentiation applied to the scalar loss.
\end{remark}

\begin{remark}[Common mistakes]
\leavevmode
\begin{enumerate}[label=(\alph*)]
  \item \emph{Forgetting the $1/m$.} If the loss is the mean, every gradient
    carries $1/m$. Omit it and the effective learning rate scales with batch
    size, which is the single most common reason a training run ``mysteriously''
    diverges when you change the batch size.
  \item \emph{Transposing the wrong factor.} Use the shape check above. Every
    \emph{weight} gradient has the form $(\text{the layer's input})^{\!\top}\bar Z$:
    $\bar W_2=R^{\!\top}\bar Z_2$ and $\bar W_1=X^{\!\top}\bar Z_1$. The bias
    gradient is instead the column sum of $\bar Z$. For the adjoint passed
    \emph{through} to the input, the transposed weight goes on the \emph{right}:
    $\bar R=\bar Z_2W_2^{\!\top}$, $\bar X=\bar ZW^{\!\top}$.
  \item \emph{Accidentally differentiating through the data.} $X$ and $Y$ are
    constants. $\bar X$ is only needed when a layer sits below; forgetting to
    stop the backward pass at the first layer is harmless, but ``updating'' $X$
    is not.
  \item \emph{Applying the mask in the wrong place.} The rule is
    $\bar Z_1=\bar R\odot M$ with $M=\mathbf{1}[Z_1>0]$ the \emph{forward} mask.
    Writing $\operatorname{relu}(\bar R)$ is wrong: that is a different,
    nonlinear operation. Multiplying $M$ into $\bar W_1$ entrywise is also wrong;
    the mask has to be applied to $\bar R$ \emph{before} the matrix product
    $X^{\!\top}(\cdot)$.
  \item \emph{Believing the gradient is the update.} The gradient is a direction
    of steepest increase; training moves \emph{against} it. Also, a zero
    gradient component means ``no first-order information'', not ``this
    parameter is useless''.
\end{enumerate}
\end{remark}

\subsection{Using the gradient: gradient descent}

\begin{definition}[Gradient descent]
With step size (learning rate) $\eta>0$, the update is
\[
  \theta^{(t+1)}=\theta^{(t)}-\eta\nabla L_D\bigl(\theta^{(t)}\bigr).
\]
\emph{Stochastic} gradient descent (SGD) replaces $L_D$ by the average loss over
a random \emph{minibatch} of examples, which is an unbiased estimator of
$\nabla L_D$ whenever the batch is sampled uniformly (see the exercises).
\end{definition}

\begin{proposition}[Descent lemma]\label{prop:descent}
Let $L$ be differentiable at $\theta$ with $\nabla L(\theta)\ne0$. Then for
every sufficiently small $\eta>0$,
$L(\theta-\eta\nabla L(\theta))<L(\theta)$.
\end{proposition}

\begin{proof}
By the definition of the total derivative,
\[
  L(\theta-\eta\nabla L(\theta))
  =L(\theta)-\eta\norm{\nabla L(\theta)}^2+o(\eta\norm{\nabla L(\theta)}).
\]
Write $r(\eta)=o(\eta\norm{\nabla L})/\eta\to0$ as $\eta\to0$. Then the
difference equals $-\eta\bigl(\norm{\nabla L}^2-r(\eta)\bigr)$, which is
negative once $\abs{r(\eta)}<\norm{\nabla L}^2$.
\end{proof}

This is a local statement only. It says a small enough step decreases the
loss; it says nothing about reaching a global minimum, and nothing about the
gap between $L_D$ and performance on unseen data. Both are real issues, taken up
in the next lesson.

\begin{example}[A complete numerical backward pass]\label{exa:numeric-backprop}
Take $m=3$ examples with $n=2$ features, a hidden layer of width $h=2$, and one
output ($k=1$):
\[
  X=\begin{pmatrix}1&2\\2&1\\1&-1\end{pmatrix},\quad
  Y=\begin{pmatrix}1\\0\\2\end{pmatrix},\quad
  W_1=\begin{pmatrix}1&-2\\2&1\end{pmatrix},\quad
  \mathbf{b}_1=\begin{pmatrix}0\\1\end{pmatrix},\quad
  W_2=\begin{pmatrix}1\\-1\end{pmatrix},\quad
  b_2=\tfrac12.
\]
\emph{Forward.} $Z_1=XW_1+\mathbf{1}\mathbf{b}_1^{\!\top}$, so
\[
  Z_1=\begin{pmatrix}5&1\\4&-2\\-1&-2\end{pmatrix},\qquad
  R=\operatorname{relu}(Z_1)=\begin{pmatrix}5&1\\4&0\\0&0\end{pmatrix},
\]
and $Z_2=RW_2+b_2\mathbf{1}$ gives
$Z_2=(4.5,\,4.5,\,0.5)^{\!\top}$. Hence
\[
  L=\frac{1}{2\cdot3}\bigl(3.5^2+4.5^2+(-1.5)^2\bigr)=\frac{34.75}{6}\approx5.792.
\]

\emph{Backward.} Applying \eqref{eq:mse-adjoint}--\eqref{eq:w1}:
\begin{align*}
  \bar Z_2&=\tfrac13(3.5,\,4.5,\,-1.5)^{\!\top}\approx(1.167,\,1.5,\,-0.5)^{\!\top},\\
  \bar W_2&=R^{\!\top}\bar Z_2
   =\begin{pmatrix}5\cdot1.167+4\cdot1.5\\ 1\cdot1.167\end{pmatrix}
   \approx\begin{pmatrix}11.833\\1.167\end{pmatrix},
  & \bar b_2&=\textstyle\sum_i\bar Z_{2,i}\approx2.167,\\
  \bar R&=\bar Z_2W_2^{\!\top}=
   \begin{pmatrix}1.167&-1.167\\1.5&-1.5\\-0.5&0.5\end{pmatrix},\\
  \bar Z_1&=\bar R\odot\mathbf{1}[Z_1>0]=
   \begin{pmatrix}1.167&-1.167\\1.5&0\\0&0\end{pmatrix},\\
  \bar W_1&=X^{\!\top}\bar Z_1
   =\begin{pmatrix}4.167&-1.167\\3.833&-2.333\end{pmatrix},
  & \bar{\mathbf b}_1&\approx\begin{pmatrix}2.667\\-1.167\end{pmatrix}.
\end{align*}
Three of the six ReLU units are off ($Z_1$ has three nonpositive entries), and
their adjoints are exactly $0$: those entries of $\bar W_1$ receive no signal at
all from this batch.
\end{example}

\begin{remark}[Reading the example as a program]
The computation is a graph, and the backward pass is the same graph reversed
with each node replaced by its adjoint rule:
\[
  X\longrightarrow Z_1\longrightarrow R\longrightarrow Z_2\longrightarrow L,
\]
with $W_1,\mathbf{b}_1$ attached to $Z_1$ and $W_2,b_2$ attached to $Z_2$. A
framework stores exactly the values needed by the mask (\eqref{eq:z1bar} needs
$Z_1$) and frees the rest, which is why activation memory dominates training
cost for large models. In code the layer is an object with two methods:
\texttt{forward(inputs) -> outputs} returning the cache, and
\texttt{backward(adjoint) -> (input\_adjoint, param\_grads)}. That is the whole
interface; \cref{ex:ml-8} asks you to justify it.
\end{remark}

\subsection{What this lesson did not do}

Two honest gaps. First, we derived gradients but not the conditions for a
minimum: the derivative test, convexity, saddle points and the fact that in
high dimension almost every stationary point is a saddle. Second, everything
here is about $L_D$, the loss on data we have. Nothing above explains why
minimising $L_D$ makes the model useful on new data, and the answer is not
calculus: it involves capacity, regularisation, and the bias--variance
decomposition.

Both come next. After that, the RL thread starts: in reinforcement learning the
loss is \emph{random}---it depends on sampled trajectories---so the object to
differentiate is an expectation, and the policy gradient theorem is exactly the
chain rule of this lesson applied to that expectation. Keep \cref{lem:adjoint-chain}
in mind; it is the only tool the policy gradient needs.
